\documentclass[10pt]{article}
\usepackage[a4paper,margin=20mm]{geometry}
\usepackage[T1]{fontenc}
\usepackage{lmodern,microtype}
\usepackage{amsmath,amssymb,mathtools,bm}
\usepackage{booktabs,tabularx,array}
\usepackage[dvipsnames]{xcolor}
\usepackage[hidelinks]{hyperref}
\usepackage{enumitem}
\setlist{nosep,leftmargin=1.5em}
\setlength{\parindent}{0pt}
\setlength{\parskip}{5pt}
\definecolor{navy}{HTML}{183B56}
\definecolor{orange}{HTML}{B65A25}
\definecolor{soft}{HTML}{FBF4EF}
\newcommand{\R}{\mathbb{R}}
\newcommand{\wrap}{\operatorname{wrap}}
\newcommand{\LN}{\operatorname{LN}}
\newcommand{\MLP}{\operatorname{MLP}}
\newcommand{\code}[1]{\texttt{#1}}
\title{\vspace{-1.4em}\color{navy}\textbf{LUMINA-2M as an AC Power-Flow Surrogate}\\
\large Exact formulation of the released HGT and PPC adaptation}
\author{Implementation note for \code{train\_valid\_test\_lumina.py}}
\date{August 2026}

\begin{document}
\maketitle
\vspace{-1.5em}
\begin{center}
\fcolorbox{orange}{soft}{\parbox{0.93\linewidth}{
\textbf{Scope.} This document describes the released LUMINA Heterogeneous Graph
Transformer loaded from \code{argonne/LUMINA-2M}, adapted to the PPC AC-PF
parquets. LUMINA was pretrained for AC-OPF-style heterogeneous inputs; here its
bus voltage output is fine-tuned or retrained against Newton PF solutions.
Optional known-operator postprocessing is excluded.
}}
\end{center}

\section{AC-PF Target}

For voltage \(\underline V_i=v_i e^{\mathrm j\theta_i}\),
\begin{equation}
 \underline{\bm S}^{\rm calc}(\bm v,\bm\theta)
 =\underline{\bm V}\odot\overline{\bm Y\underline{\bm V}}.
\end{equation}
The mismatch is
\begin{equation}
 \Delta\bm P=\bm P^{\rm set}-\Re\underline{\bm S}^{\rm calc},\qquad
 \Delta\bm Q=\bm Q^{\rm set}-\Im\underline{\bm S}^{\rm calc}.
\end{equation}
The adapted LUMINA model learns
\begin{equation}
 f_\Theta(\mathcal G,\underline{\bm S}^{\rm set})
 \longmapsto(\hat{\bm v},\hat{\bm\theta})
 \approx(\bm v^\star,\bm\theta^\star).
\end{equation}

\section{Adapted Heterogeneous Graph}

LUMINA uses four node types:
\begin{equation}
 \mathcal V=\mathcal V_b\cup\widetilde{\mathcal V}_g\cup
 \mathcal V_d\cup\mathcal V_s,
\end{equation}
representing buses, proxy generators, loads, and shunts. It uses eight directed
relation types:
\begin{align}
 &b\xrightarrow{\rm ac\ line}b,\qquad
 b\xrightarrow{\rm transformer}b,\\
 &\tilde g\mathrel{\overset{\rm generator\ link}{\rightleftarrows}}b,\qquad
 d\mathrel{\overset{\rm load\ link}{\rightleftarrows}}b,\qquad
 s\mathrel{\overset{\rm shunt\ link}{\rightleftarrows}}b.
\end{align}
The source LUMINA schema supports physical generator nodes. Since the PPC PF
parquet is bus-aggregated, the adapter constructs one proxy generator on every
PV/REF bus.

\subsection{Node features}

The bus row is
\begin{equation}
 \bm x_i^b=[V_{n,i}^{\rm kV},v_i^{\min},v_i^{\max},
 \mathbb 1_{\rm PQ},\mathbb 1_{\rm PV},\mathbb 1_{\rm REF},
 \mathbb 1_{\rm ISO}]\in\R^7.
\end{equation}
The generator, load, and shunt rows are
\begin{align}
 \bm x_{\tilde g_i}^g&=[m_{\rm base},P_g,P_g^{\min},P_g^{\max},Q_g,
 Q_g^{\min},Q_g^{\max},v_g^{\rm set},c_2,c_1,c_0]\in\R^{11},\\
 \bm x_{d_i}^d&=[P_{d,i},Q_{d,i}]\in\R^2,\\
 \bm x_{s_i}^s&=[B_i^{\rm sh},G_i^{\rm sh}]\in\R^2.
\end{align}
In the PF adaptation,
\begin{itemize}
 \item \(P_g,Q_g\) may be filled with the known bus-level injection
       (\code{gen\_setpoint\_mode=known});
 \item generator limits are synthesized around that injection and costs are zero;
 \item \(v_g^{\rm set}=v_i^0\) at proxy-generator buses;
 \item load channels are the clipped negative net injection;
 \item the unspecified \(Q\) slot at PV/REF buses is zeroed before graph construction.
\end{itemize}

The bus row itself does \emph{not} include \((v_i^0,\theta_i^0)\). Only the
generator setpoint channel exposes \(v_i^0\) at PV/REF buses. Therefore the
manual-flat versus DC-start angle is not directly provided to released LUMINA.

\subsection{Branch attributes supplied by the adapter}

AC-line and transformer rows follow the PGLib/OPF schema:
\begin{align}
 \bm e_{ij}^{\ell}&=[\theta^{\min},\theta^{\max},b_{fr},b_{to},r,x,
 S_A,S_B,S_C]\in\R^9,\\
 \bm e_{ij}^{\tau}&=[\theta^{\min},\theta^{\max},r,x,S_A,S_B,S_C,
 \tau,\phi,b_{fr},b_{to}]\in\R^{11}.
\end{align}
They are reconstructed from PPC branch admittance, tap, phase shift, and shunts.

\paragraph{Implementation-critical fact.}
The released \code{HGT.forward} accepts an \code{edge\_attr\_dict} argument but
does not pass it into \code{HGTConv}. Consequently, the present LUMINA model uses
branch attributes only indirectly through graph construction/classification;
the numerical values \((r,x,b,\tau,\phi)\) do not enter message attention. The
relations distinguish ``AC line'' from ``transformer'', but two lines with the
same endpoints and different impedance are indistinguishable to the HGT core.

\section{Type Projections}

Each node type \(t\in\{b,g,d,s\}\) has its own input projection:
\begin{equation}
 \bm h_i^{t,(0)}=\operatorname{ReLU}(W_{\rm in}^{t}\bm x_i^t+\bm b_{\rm in}^{t}).
\end{equation}
The released checkpoint configuration uses
\begin{equation}
 d=128,\qquad L=6,\qquad H=2,\qquad p_{\rm drop}=0.2439835578.
\end{equation}

\section{Heterogeneous Graph Transformer Layers}

Let relation \(r=(s,\varrho,t)\) connect source type \(s\) to target type \(t\).
For head \(h\), HGT forms type-specific queries, keys, and values:
\begin{equation}
 \bm q_i^{t,h}=W_Q^{t,h}\bm h_i^t,\qquad
 \bm k_j^{s,h}=W_K^{s,h}\bm h_j^s,\qquad
 \bm v_j^{s,h}=W_V^{s,h}\bm h_j^s.
\end{equation}
The relation transforms key/value information and supplies a learned relation
prior. An operator-level form is
\begin{align}
 e_{ji}^{r,h}&=\frac{(\bm q_i^{t,h})^\top A_r^h\bm k_j^{s,h}}
 {\sqrt{d_h}}+\mu_r^h,\\
 \alpha_{ji}^{r,h}&=\operatorname{softmax}_{(r',u):u\to i}(e_{ui}^{r',h}),\\
 \bm m_i^{t,h}&=\sum_{r=(s,\cdot,t)}\sum_{j\in\mathcal N_r(i)}
 \alpha_{ji}^{r,h}M_r^h\bm v_j^{s,h}.
\end{align}
Thus different relations learn different transformations even when they share a
target node type. The heads are combined, transformed by a target-type output
map, and mixed with a type-specific skip path:
\begin{equation}
 \bm h_i^{t,(\ell+1)}=\operatorname{ReLU}
 \left(\operatorname{HGTConv}^{(\ell)}_t
 (\{\bm h^{(\ell)}\},\{\mathcal E_r\})_i\right).
\end{equation}
Dropout follows the input projection and every HGT layer during training.

Unlike GraphKit's HGNS, LUMINA has no layerwise voltage decoding, no branch-flow
calculation, and no AC residual feedback inside its forward pass.

\section{Terminal Output Heads}

Only after the sixth HGT layer, two type-specific linear heads are applied:
\begin{align}
 [z_{\theta,i},z_{v,i}]&=W_{\rm out}^{b}\bm h_i^{b,(L)}+\bm b_{\rm out}^{b},\\
 [z_{P,g},z_{Q,g}]&=W_{\rm out}^{g}\bm h_g^{g,(L)}+\bm b_{\rm out}^{g}.
\end{align}
With min--max scaling enabled, the magnitude and generator outputs are bounded:
\begin{align}
 \hat v_i&=v_i^{\min}+(v_i^{\max}-v_i^{\min})\sigma(z_{v,i}),\\
 \hat P_g&=P_g^{\min}+(P_g^{\max}-P_g^{\min})\sigma(z_{P,g}),\\
 \hat Q_g&=Q_g^{\min}+(Q_g^{\max}-Q_g^{\min})\sigma(z_{Q,g}).
\end{align}
The angle is not bounded by min--max scaling; the adapter wraps it:
\begin{equation}
 \hat\theta_i=\operatorname{atan2}(\sin z_{\theta,i},\cos z_{\theta,i}).
\end{equation}
The AC-PF path consumes
\begin{equation}
 \widehat{\bm V}_{\rm state}=\{(\hat v_i,\hat\theta_i)\}_{i=1}^{N}.
\end{equation}
The generator output head is unused under the default voltage-only PF objective.
The forward function also does not pin known PF quantities: it predicts
magnitude at PV/REF buses and angle at the REF bus, and supervision must learn
those identities.

\section{External AC-PF Physics Loss}

The supervised term is
\begin{equation}
 \mathcal L_{\rm state}=\frac1N\sum_i
 \left[(\hat v_i-v_i^\star)^2+
 \wrap(\hat\theta_i-\theta_i^\star)^2\right].
\end{equation}
Let
\begin{equation}
 \mathcal M_P=\mathcal V_b\setminus\mathcal V_{\rm REF},\qquad
 \mathcal M_Q=\mathcal V_{\rm PQ}.
\end{equation}
The external physics term uses the exact PPC \(\bm Y\), independent of the HGT
edge representation:
\begin{equation}
 \mathcal L_{\rm phys}=\frac1{|\mathcal M_P|+|\mathcal M_Q|}
 \left[\sum_{i\in\mathcal M_P}\rho(\Delta P_i)+
 \sum_{i\in\mathcal M_Q}\rho(\Delta Q_i)\right],
\end{equation}
with default
\begin{equation}
 \rho(r)=\log\cosh(\operatorname{clip}(r,-30,30)).
\end{equation}
The standard objective is
\begin{equation}
 \boxed{\mathcal L=\mathcal L_{\rm state}+10^{-2}\mathcal L_{\rm phys}.}
\end{equation}
Therefore LUMINA can still learn impedance-sensitive AC behavior through the
external loss gradient, even though impedance is absent from its message
features. It must infer much of that behavior from topology, relation type,
nominal voltage, injections, and repeated-grid training examples.

Training physics is evaluated through complex64 casts in the script; reported
validation/test residual metrics use the complex128 dataset/Ybus path.

\section{Interpretive Summary}
\begingroup\small
\begin{tabularx}{\linewidth}{>{\bfseries}p{0.25\linewidth}X}
\toprule
Property & LUMINA PF adaptation \\
\midrule
Graph & Bus, proxy-generator, load, and shunt nodes with eight directed relation types.\\
Electrical edge encoding & Line/transformer attributes are constructed, but the released HGT forward currently ignores numerical edge attributes.\\
Neural core & Six HGTConv layers, hidden width 128, two heads, and dropout about 0.244.\\
Voltage decoder & One terminal bus head; magnitude sigmoid-scaled to the voltage band and angle wrapped.\\
PF constraints inside forward & No PQ/PV/REF pinning, no branch-flow operator, and no residual feedback.\\
Physics outside forward & Exact-Ybus robust AC-PF loss and complex128 evaluation metrics.\\
Start-state use & Start magnitude appears only as proxy-generator \(v_g^{\rm set}\); start angle is omitted.\\
Pretraining & Released LUMINA-2M weights may be loaded and fine-tuned.\\
\bottomrule
\end{tabularx}

\section*{Code provenance}
The formulation follows \code{lumina\_inference/model/hetero\_model.py}, the
released \code{lumina\_config.json}, and the AC-PF adapter in
\code{train\_valid\_test\_lumina.py}, as present in the local workspace in
August 2026.
\endgroup

\end{document}
